\section{Round-sixteen positive closure: state-dependent Nisio resolvents and a diagonal hierarchy limit}
\label{sec:r16-b4}

The exact finite system, its Koopman observables, push-forward laws, factorial
hierarchies, and limiting density value are kept distinct.  Comparison is not
proved by applying one contact flow to two different density states.  It is
derived from the dynamic action, its exact concatenation law, and a nonlinear
resolvent identity.

\subsection{Exact microscopic objects}

Let \(\Phi_t^\varepsilon\) be the deterministic hard-sphere flow on the
Liouville phase space and let
\[
 U_t^\varepsilon F=F\circ\Phi_t^\varepsilon,
 \qquad
 (\Phi_t^\varepsilon)_\#\rho
\]
be, respectively, the Koopman semigroup on observables and the push-forward
semigroup on probability laws.  On the exponential-moment law class
\[
 \mathcal P_\alpha=
 \left\{\rho:\int e^{\alpha N+\beta E_N}d\rho<\infty\right\},
\]
the factorial correlations \(G(\rho)=(g_k)_{k\ge0}\) satisfy
\[
 \sum_{k\ge0}{e^{\alpha'k}\over k!}
 \|g_k\|_{L^1_\beta}<\infty,
 \qquad \alpha'<\alpha.
\]

\begin{theorem}[Law--hierarchy determinacy and intertwining]
\label{thm:r16-b4-intertwiner}
The map \(G:\mathcal P_\alpha\to\mathfrak X_{\alpha'}\) is injective, its
image is closed under the short-time BBGKY evolution, and
\[
 G((\Phi_t^\varepsilon)_\#\rho)
 =\mathscr U_t^\varepsilon G(\rho).
\]
For every bounded exponential cylinder, its expectation is the absolutely
convergent factorial-moment series determined by \(G(\rho)\).
\end{theorem}

\begin{proof}
The exponential particle-number moment makes the probability generating
functional
\[
 \mathcal Z_\rho(h)=
 \mathbb E_\rho\prod_i(1+h(z_i))
 =\sum_{k\ge0}{1\over k!}\int h^{\otimes k}g_k
\]
analytic on an \(L^\infty\) ball.  Equality of all correlations therefore
implies equality of the generating functional and of the law.  Liouville
transport differentiates this series term by term on a smaller exponential
weight and gives the BBGKY hierarchy, proving the intertwining.  The same
normal convergence evaluates every bounded exponential cylinder.
\end{proof}

The exponential collision Hamiltonian is not the exponential conjugate of the
interior Koopman derivation.  It comes from the ensemble boundary jump.

\begin{theorem}[Ensemble boundary generator]
\label{thm:r16-b4-boundary}
For a smooth density cylinder \(F\), \(p=\delta F/\delta f\), and bounded
actual-contact work \(\psi\), the exact finite ensemble log-Laplace derivative
has boundary contribution
\[
 {1\over2}\int ff_*B
 \bigl(e^{\Delta p+\psi}-1\bigr)
\]
after the Boltzmann--Grad limit.  Together with free transport this is the B2
Hamiltonian \(\mathbb H(f,p,\psi)\).
\end{theorem}

\begin{proof}
Apply the Green formula of B2 to the exponential cylinder before taking the
logarithm.  At an actual incoming contact the ratio of post- to pre-collision
exponential weights is \(e^{\Delta p+\psi}\).  The B2 all-contact source
theorem controls recollisions and correlations and passes the complete tilted
boundary flux to \(ff_*B/2\).  The interior derivation contributes only
\(\langle f,v\cdot\nabla p\rangle\).
\end{proof}

\subsection{The dynamic action and compact state space}

Let
\[
 \mathcal E_M=
 \left\{f\ge0:\langle1+|v|^m+e^{\alpha|v|^2},f\rangle\le M\right\}
\]
with the weighted weak topology generated by a countable dense set of tests.
For a balanced pair \((f,\Gamma)\) on \([s,t]\), define the transition action
\[
 \mathcal A_{s,t}(f,\Gamma)
 =\int_s^t\ell(d\Gamma/dA_f)dA_f.
\]
No initial preparation cost is included in a transition action.

\begin{lemma}[Action compactness and concatenation]
\label{lem:r16-b4-compact}
For fixed endpoints in a compact regular initial set and bounded action,
balanced paths are compact in the weighted weak Skorokhod topology.  If two
balanced paths meet at time \(r\), their concatenation is balanced and
\[
 \mathcal A_{s,t}=\mathcal A_{s,r}+\mathcal A_{r,t}.
\]
Conversely restriction to subintervals preserves the action identity.
\end{lemma}

\begin{proof}
Entropy coercivity controls contact mass and its weighted tails.  The weak
balance then controls time increments against every weighted smooth test.
Energy conservation and the exponential velocity Lyapunov estimate give
uniform tightness in velocity.  Prokhorov and the Jakubowski modulus criterion
give compactness.  Balance is linear and closed; splitting and joining the
contact measure at \(r\) proves concatenation.  The entropy integral is local
in time, giving exact additivity.
\end{proof}

For bounded uniformly continuous \(\Phi\), define
\[
 (S_t\Phi)(f_0)=
 \sup_{(f,\Gamma):f(0)=f_0}
 \{\Phi(f_t)-\mathcal A_{0,t}(f,\Gamma)\}.
\]

\begin{theorem}[Nisio semigroup]
\label{thm:r16-b4-nisio}
The value is finite and attained after relaxation, maps
\(\operatorname{BUC}(\mathcal E_M)\) into itself, is monotone, cash additive,
and a sup-norm contraction.  It satisfies
\[
 S_{t+r}=S_tS_r.
\]
The continuity modulus is obtained by state-dependent recovery paths; no
contact flow is reused from a different initial state.
\end{theorem}

\begin{proof}
Compactness and lower semicontinuity from
Lemma~\ref{lem:r16-b4-compact} give attainment.  Concatenating admissible paths
proves \(S_{t+r}\ge S_tS_r\); splitting a path at time \(r\) proves the reverse
inequality.  Monotonicity, cash additivity, and contraction are immediate from
the terminal payoff.

For continuity in \(f_0\), approximate a nearly optimal path by a smooth
strictly positive exposed B2 path.  Solve the controlled Boltzmann equation
with the same bounded source but the perturbed initial state.  Short-time
well-posedness and Gronwall stability give a nearby balanced path and action
change tending to zero.  This transports the source, not the original contact
measure, and is admissible from each state separately.  Approximation and the
reverse argument give the BUC modulus.
\end{proof}

\subsection{Nonlinear resolvent and full range}

For \(\lambda>0\), define the discounted value
\[
 (R_\lambda h)(f_0)=
 \sup_{(f,\Gamma):f(0)=f_0}
 \int_0^\infty e^{-\lambda t}
 \left[\lambda h(f_t)-
 {d\mathcal A\over dt}(f_t,\Gamma_t)\right]dt.
\]

\begin{theorem}[Nonlinear resolvent identity and \(m\)-dissipativity]
\label{thm:r16-b4-resolvent}
For all \(\lambda,\mu>0\),
\[
 R_\lambda h
 =R_\mu\left({\mu\over\lambda}h+
   \left(1-{\mu\over\lambda}\right)R_\lambda h\right)
\]
with the standard equivalent form when \(\mu>\lambda\).  Each
\(R_\lambda\) is a sup-norm contraction and its range is the whole BUC state
space.  The graph
\[
 \mathscr H=
 \{(R_\lambda h,\lambda(R_\lambda h-h)):h\in\operatorname{BUC}\}
\]
is independent of \(\lambda\), closed, dissipative, and has full resolvent
range.  It is therefore \(m\)-dissipative, and its Crandall--Liggett semigroup
is \(S_t\).
\end{theorem}

\begin{proof}
Split the discounted integral at a deterministic time and apply the dynamic
programming identity of Theorem~\ref{thm:r16-b4-nisio}.  Integrating the split
identity against either exponential clock and using Fubini gives the nonlinear
resolvent identity.  The same variational formula gives contraction.  Since
\(R_\lambda h\) is defined for every bounded uniformly continuous \(h\), the
range condition \((I-\lambda^{-1}\mathscr H)u=h\) is exact.  The resolvent
identity makes the graph independent of \(\lambda\); contraction gives
dissipativity and closedness.  The nonlinear Hille--Yosida theorem identifies
its semigroup, while uniqueness of the dynamic programming semigroup gives
\(S_t\).
\end{proof}

\subsection{Graph-core generator and comparison}

Let \(\mathscr D\) be the smooth density cylinders whose variational derivative
has bounded collision increment and weighted transport derivative.  On this
core define \(\mathbb H\) by
Theorem~\ref{thm:r16-b4-boundary}.

\begin{theorem}[Primal resolvent comparison]
\label{thm:r16-b4-comparison}
The closure of \((\mathbb H,\mathscr D)\) in the BUC graph norm is
\(\mathscr H\).  Hence the discounted Hamilton--Jacobi equation
\[
 \lambda u-\mathbb H(f,Du)=\lambda h
\]
has the unique BUC viscosity/resolvent solution \(u=R_\lambda h\), and the
finite-horizon equation has the unique value \(S_{T-t}\Phi\).
\end{theorem}

\begin{proof}
For a core cylinder, restrict the action value to a short interval.  Fenchel
duality gives the upper generator inequality, while the B2 bounded exposing
source gives a state-dependent recovery path and the matching lower inequality.
Thus the core graph lies in \(\mathscr H\).  Conversely, Yosida approximants
\(u_n=R_n h\) lie in the resolvent graph.  Smooth their finite determining
coordinates, truncate velocity with an exponential Lyapunov corrector, and
use the B2 positive entropic projection to restore balance.  This produces
core cylinders converging in graph norm.  Equality of closures follows.
Uniqueness is then the contraction comparison of the \(m\)-dissipative
resolvent, not a doubled-variable argument with uncontrolled state-dependent
controls.
\end{proof}

\subsection{A single diagonal hierarchy corrector}

For a core cylinder \(F\), solve the finite-label triangular corrector equations
in the BBGKY graph domain through level \(K\).  The \(j\)-label term carries
weight \(1/j!\).

\begin{theorem}[Diagonal perturbed tests]
\label{thm:r16-b4-corrector}
There are \(K(\varepsilon)\uparrow\infty\) and exact hierarchy cylinders
\(F_\varepsilon\) such that
\[
 F_\varepsilon\to F,
 \qquad
 \mathscr H_\varepsilon F_\varepsilon
 \to\mathbb H(F)
\]
locally uniformly on every exponential-moment action sublevel.  The sequence
has a uniform graph bound and is compatible with compact containment.
\end{theorem}

\begin{proof}
At fixed \(K\), solve the triangular terminal-value Duhamel equations on the
observable graph domain; B2 controls each connected \(j\)-label defect by
\(j!C^jT^{j-1}\).  The factorial hierarchy norm therefore makes the corrected
series tail at most \(C\rho^K\).  Recollision defects are at most
\(C_K\varepsilon^{\alpha_K}\).  Choose \(K(\varepsilon)\) increasing so slowly
that both \(\rho^K\) and \(C_K\varepsilon^{\alpha_K}\) tend to zero.  The same
bounds control the graph norm and exponential moments.  This gives one
recovery sequence rather than two iterated limits.
\end{proof}

\begin{theorem}[Microscopic nonlinear semigroup convergence]
\label{thm:r16-b4-limit}
For exact initial laws in a compact exponential-moment class, the centered
finite-volume ensemble log-Laplace values converge locally uniformly to
\(S_t\).  The normalizers are the exact finite-volume pressures.  Under the B3
central scaling, the values converge jointly to the risk-sensitive Gaussian
semigroup with the complete B1 initial and B3 density/contact covariance.
\end{theorem}

\begin{proof}
The diagonal corrector gives upper and lower nonlinear generator convergence
on a separating graph core.  Exponential compact containment follows from B2.
The nonlinear Trotter--Kato theorem for \(m\)-dissipative graphs and
Theorem~\ref{thm:r16-b4-comparison} identify the limit with \(S_t\).
Centering by the exact finite pressure makes each likelihood ratio mean one.
B3 gives joint process Gaussianity and uniform \(L^{1+\delta}\) bounds on a
larger source ball; Taylor expansion of the exact normalizer yields the
quadratic risk-sensitive generator.
\end{proof}

\begin{theorem}[Round-sixteen B4 closure]
\label{thm:r16-b4-main}
The microscopic law and hierarchy are exactly intertwined on an
determinate exponential-moment class; the ensemble boundary generates the
Boltzmann Hamiltonian; the dynamic action constructs a genuine Nisio semigroup
and \(m\)-dissipative resolvent; a single diagonal hierarchy corrector gives
microscopic convergence; and comparison uses no inadmissible common contact
flow.
\end{theorem}

\begin{proof}
Combine Theorems~\ref{thm:r16-b4-intertwiner},
\ref{thm:r16-b4-boundary}, \ref{thm:r16-b4-nisio},
\ref{thm:r16-b4-resolvent}, \ref{thm:r16-b4-comparison},
\ref{thm:r16-b4-corrector}, and \ref{thm:r16-b4-limit}, and
Lemma~\ref{lem:r16-b4-compact}.
\end{proof}
